\documentclass[11pt]{article}
\usepackage{fontspec}
\usepackage{xeCJK}
\setmonofont{DejaVuSansMono.ttf}[BoldFont=DejaVuSansMono-Bold.ttf,Scale=0.82]
\usepackage{xurl}
\usepackage[a4paper,margin=25mm]{geometry}
\usepackage{amsmath,amssymb,amsthm,mathtools}
\usepackage{xcolor}
\usepackage{booktabs,array,enumitem,fancyhdr,hyperref}
\hypersetup{colorlinks=true,linkcolor=blue!50!black,urlcolor=blue!50!black}
\setlength{\emergencystretch}{2em}
\setlength{\parskip}{0.3em}
\linespread{1.16}
\allowdisplaybreaks
\newtheorem{theorem}{定理}[section]
\newtheorem{proposition}[theorem]{命题}
\newtheorem{lemma}[theorem]{引理}
\theoremstyle{definition}
\newtheorem{definition}[theorem]{定义}
\theoremstyle{remark}
\newtheorem{remark}[theorem]{说明}
\newcommand{\dd}{\,\mathrm d}
\newcommand{\ii}{\mathrm i}
\newcommand{\Rea}{\operatorname{Re}}
\newcommand{\Ima}{\operatorname{Im}}
\newcommand{\Li}{\operatorname{Li}}
\newcommand{\supp}{\operatorname{supp}}
\pagestyle{fancy}
\fancyhf{}
\fancyhead[L]{\small 学习笔记}
\fancyhead[R]{\small 初稿}
\fancyfoot[C]{\thepage}
\setlength{\headheight}{15pt}
\author{}
\date{2026年10月22日初稿}
\makeatletter
\renewcommand{\maketitle}{%
  \thispagestyle{plain}%
  \begin{center}
    {\LARGE\bfseries \@title\par}
    \vspace{0.6em}
    {\small \@date\par}
  \end{center}
  \vspace{0.6em}}
\makeatother
\newcommand{\draftnotice}{\noindent{\small
本文为供逐篇审阅的学习笔记初稿，不是新增研究成果。
文中区分直接证明与引用结果；正文不依赖本文。
本轮已作公式、条件及排版初步核查，尚待作者逐句审定。}\par\medskip}

\usepackage{fancyvrb}
\newfontfamily\symfont{DejaVuSans.ttf}[Scale=0.82]
\fvset{fontsize=\small,xleftmargin=1.5em,frame=leftline,framesep=0.8em}
\XeTeXlinebreaklocale "zh"
\renewcommand{\proofname}{证明}
\date{2026年10月26日初稿}
\title{AI 时代的形式化验证：以 $\zeta$ 函数为例}
\begin{document}
\maketitle

\noindent{\small 本篇以 Lean 4 及其数学库 mathlib 中关于 $\zeta$ 函数的定理为例，说明“经过形式化验证”保证什么、不保证什么，
以及在 AI 大量生成证明的情况下应当怎样阅读这类声明。文中的 Lean 代码取自本仓库 \texttt{Lean示例/}，均已由 Lean 检查；
形式化验证的一般原理见学习笔记第 16 篇，此处不重复。所涉近期事件截至 2026 年 10 月下旬，均注明出处。}\par\medskip

\section{问题}
传统上，一个证明是否可信，取决于同行能否读懂并逐步核对。当证明由 AI 生成、篇幅达到数百页、而人类读者有限时，
这一方式遇到了瓶颈。形式化验证提供另一条途径：把命题和证明写成证明助手能够机械检查的形式，
由一个很小的程序（内核）判断每一步推理是否合法。

近来越来越多的数学声明附有“已在 Lean 中形式化”的说明。本篇讨论三个问题：
\begin{enumerate}
 \item 一条形式化陈述怎样读，它与通常的数学陈述有哪些差别；
 \item Lean 接受一个证明，究竟保证了什么；
 \item 证明者不可信（例如未经审查的 AI）时，还需要哪些额外的检查。
\end{enumerate}

\section{一条形式化陈述怎样读}
mathlib 中黎曼猜想的定义如下：
\begin{Verbatim}
def RiemannHypothesis : Prop :=
  ∀ (s : ℂ) (_ : riemannZeta s = 0) (_ : ¬∃ n : ℕ, s = -2 * (n + 1))
    (_ : s ≠ 1), s.re = 1 / 2
\end{Verbatim}
\texttt{Prop} 表示这是一个命题，而不是已经证明的定理。逐项读出：对一切复数 $s$，若 $\zeta(s)=0$，
$s$ 不是 $-2,-4,-6,\dots$ 中的某一个，且 $s\ne1$，则 $\Rea s=\tfrac12$。

这条陈述有三处值得注意。
\begin{enumerate}
 \item \textbf{类型。}\texttt{s : ℂ} 指明 $s$ 是复数；\texttt{n : ℕ} 指明 $n$ 是自然数（含 $0$），
 因此 \texttt{-2 * (n + 1)} 恰好取遍 $-2,-4,\dots$。若误写成 \texttt{-2 * n}，就会把 $s=0$ 也排除在外（因 $\zeta(0)=-\tfrac12\ne0$，这碰巧不改变命题的真假）。
 形式化陈述中的每个符号都有确定的类型，这消除了自然语言中的许多歧义，也带来了新的出错方式。
 \item \textbf{定义的依赖。}陈述中的 \texttt{riemannZeta} 是 mathlib 中的一个定义，经由 Hurwitz $\zeta$ 函数、
 Mellin 变换与 $\theta$ 函数逐层构造，最终依赖实数与复数的构造。这条陈述是否“就是”黎曼猜想，
 取决于这一长串定义是否与数学上的 $\zeta$ 一致。mathlib 为此证明了大量性质（$\Rea s>1$ 时等于 $\sum n^{-s}$、
 函数方程、Euler 乘积等），作为“定义正确”的旁证。
 \item \textbf{多余的条件。}条件 $s\ne1$ 在数学上多余，因为 $\zeta$ 在 $s=1$ 无定义。在 mathlib 中它同样多余（见下节），
 但写上它，陈述就不依赖于 $s=1$ 处的约定。
\end{enumerate}

\section{处处有定义的函数}
\subsection{$\zeta(1)$ 的约定值}
Lean 中的函数必须在定义域的每一点取值，\texttt{riemannZeta} 的定义域是整个 $\mathbb C$，因此 \texttt{riemannZeta 1} 必须是某个复数。
mathlib 的构造方式（其中除以零的结果规定为零）使
\begin{equation}
 \texttt{riemannZeta 1}=\frac{\gamma-\log(4\pi)}{2}\approx-0.9769 .\tag{3.1}
\end{equation}
\begin{proposition}
设 $\Lambda(s)=\pi^{-s/2}\Gamma(s/2)\zeta(s)$ 为完备 $\zeta$ 函数，则它在 $s=1$ 处 Laurent 展开的常数项等于 $(\gamma-\log4\pi)/2$。
\end{proposition}
\begin{proof}
在 $s=1$ 附近，$\zeta(s)=\dfrac1{s-1}+\gamma+O(s-1)$。又 $g(s)=\pi^{-s/2}\Gamma(s/2)$ 在 $s=1$ 处解析，$g(1)=\pi^{-1/2}\Gamma(\tfrac12)=1$，
\[
 \frac{g'(1)}{g(1)}=-\frac12\log\pi+\frac12\,\psi\!\left(\tfrac12\right)=-\frac12\log\pi-\frac\gamma2-\log2 ,
\]
其中用到 $\psi(\tfrac12)=-\gamma-2\log2$。于是 $g(s)=1+g'(1)(s-1)+O\bigl((s-1)^2\bigr)$，
\[
 \Lambda(s)=\frac1{s-1}+\gamma+g'(1)+O(s-1)=\frac1{s-1}+\frac{\gamma-\log4\pi}{2}+O(s-1). \qedhere
\]
\end{proof}
这个约定值并非任意：它是 $\Lambda$ 去掉极点部分后在 $s=1$ 处的值，也等于 $\lambda_1-1$，
其中 $\lambda_1=1+\tfrac\gamma2-\tfrac12\log4\pi$ 是 Li 判据的第一个系数（学习笔记第 10 篇）。
但它不是 $\zeta(s)-\dfrac1{s-1}$ 在 $s=1$ 处的极限 $\gamma$。两种“自然”的补充定义给出不同的值，正说明约定值不携带关于极点的数学信息。

由此，mathlib 的定理
\begin{Verbatim}[commandchars=\\\{\}]
theorem riemannZeta_ne_zero_of_one_le_re {\symfont ⦃}s : ℂ{\symfont ⦄} (hs : 1 ≤ s.re) :
    riemannZeta s ≠ 0
\end{Verbatim}
在 $s=1$ 处断言的是 $(\gamma-\log4\pi)/2\ne0$，与 $\zeta$ 在 $s=1$ 有极点这一事实无关。它在 $s\ne1$ 处的内容才是本书第 9 章的定理。

\subsection{$\Gamma$ 在极点处的值与函数方程的条件}
同样，mathlib 规定 \texttt{Complex.Gamma} 在 $0,-1,-2,\dots$ 处取 $0$。函数方程在 mathlib 中写作
\begin{Verbatim}
theorem riemannZeta_one_sub {s : ℂ} (hs : ∀ n : ℕ, s ≠ -n) (hs' : s ≠ 1) :
    riemannZeta (1 - s) =
      2 * (2 * π) ^ (-s) * Gamma s * cos (π * s / 2) * riemannZeta s
\end{Verbatim}
条件 \texttt{hs} 排除 $s=0,-1,-2,\dots$。它在 Lean 中是必需的，而不只是习惯：取 $s=0$，左边是约定值 (3.1)，不为零；
右边含因子 \texttt{Gamma 0}，按约定为 $0$，于是右边为 $0$。去掉 \texttt{hs}，这条陈述在 Lean 中是假的。
数学上 $s=0$ 处右边是“极点乘零点”的形式，公式只能通过取极限理解；Lean 不做这种理解，条件必须写明。

这类约定是形式化数学的常规做法：让函数处处有定义，代价是在每条定理中写清适用范围。
读者遇到“mathlib 中某公式对一切 $s$ 成立”时，应先查明它在例外点上依赖哪些约定。

\section{Lean 接受一个证明，保证了什么}
\subsection{保证的部分}
Lean 接受一个定理，意味着内核逐步检查了一个证明项，确认它在 Lean 的逻辑中推出了所写的陈述。
可信度依赖于：内核与逻辑的可靠性；所用的公理；所依赖的库定义（见学习笔记第 16 篇）。
其中公理可以直接查询：命令 \texttt{\#print axioms} 列出一个定理直接或间接依赖的全部公理。
标准情形只有 \texttt{propext}、\texttt{Classical.choice}、\texttt{Quot.sound} 三条；
若出现 \texttt{sorryAx}，说明证明中有未完成的部分；若出现其他名字，则是另外引入的公理，定理只在该公理成立时才成立。

\subsection{不保证的部分：陈述}
Lean 不判断陈述是否就是想要的命题。下面这条定理被 Lean 接受（本仓库 \texttt{Lean示例/\allowbreak LeanExamples/\allowbreak Ch11\_RiemannHypothesis.lean}）：
\begin{Verbatim}
example (s : ℂ) (h1 : 1 < s.re) (h2 : riemannZeta s = 0) :
    s.re = 1 / 2 :=
  absurd h2 (riemannZeta_ne_zero_of_one_le_re h1.le)
\end{Verbatim}
它的字面意思是“若 $\Rea s>1$ 且 $\zeta(s)=0$，则 $\Rea s=\tfrac12$”，结论与 RH 相同，证明只有一行。
原因是两个假设不能同时成立，而从矛盾的假设可以推出任何结论。这是形式化中最典型的一类陷阱：
\emph{假设过强}，以致定理空洞成立。与之相对的是\emph{结论过弱}，或者把定义写错，使陈述说的是另一件事。

因此，一条“已经形式化的定理”是否有数学价值，要由人回答两个问题：形式化陈述是否忠实于原意；它所依赖的定义是否正确。
Lean 把需要人工审查的部分从整个证明缩小到陈述与定义，但没有消除这一部分。

\section{证明者不可信的情形}
Lean 的官方参考手册在“验证证明”一章中区分两种情形：\emph{诚实}的证明尝试可能有错误，但不试图欺骗系统；
\emph{恶意}的尝试则可能利用系统漏洞，或用元程序在检查之前做手脚。手册明确把未经审查的 AI 生成证明归入后一类，
并给出由弱到强的几层检查：
\begin{enumerate}
 \item 编辑器显示通过，或 \texttt{lake build} 无错误。这只说明当前文件被接受，不能发现依赖中的 \texttt{sorry}。
 \item \texttt{\#print axioms}，拒绝 \texttt{sorryAx} 与非标准公理。
 \item 用 \texttt{lean4checker} 把编译结果重新送入内核检查，排除某些绕过内核的漏洞。
 \item 最强的做法是 \texttt{comparator}：由可信的一方单独写出\emph{挑战文件}，其中只有定理陈述而没有证明；
 待检查的证明在隔离环境中编译并导出，再在隔离环境之外由 Lean 内核或独立实现的外部检查器重新检查，
 并逐字比对它证明的陈述是否与挑战文件一致。
\end{enumerate}
第 4 层的要点在于分工：证明可以来自任何人或任何程序，而陈述由可信的人写定。审查者不必阅读证明，
只需审查挑战文件中的陈述及其所用定义是否正确。这正对应上一节所说的“需要人工审查的部分”。

Tao 在评论 AI 完成形式化项目时也指出，项目的组织者今后需要把全部目标明确写出，因为 AI 只会优化被明确陈述的目标。
形式化把“证明是否正确”交给机器，同时把“陈述是否正确”的责任更集中地留给人。

\section{两个实例}
\subsection{强形式素数定理（2025）}
2024 年 1 月，Tao 与 Kontorovich 发起在 Lean 中形式化素数定理（带经典误差项）及相关结果的项目 PrimeNumberTheoremAnd。
到 2025 年中，项目取得部分进展，主要困难在于复分析基础结果的缺失。2025 年 9 月，Math Inc.\ 公司报道其自动形式化代理 Gauss
用约三周完成了强形式素数定理的形式化，代码约 $25\,000$ 行，期间补齐了所需的复分析结果。
（在此之前，素数定理已于 2007 年在 Isabelle 中、2009 年在 HOL Light 中形式化，分别采用初等证明与解析证明。）

这一例子说明，AI 在形式化中的作用之一是承担大量“已知数学”的翻译工作。陈述本身（素数定理）没有争议，验证的价值在于证明的完整性。

\subsection{“拟黎曼猜想”（2026）}
2026 年 10 月 6 日，OpenAI 公开一批由其内部模型生成的数学稿件（据报道共 722 份），其中据报道约四到六成附有 Lean 证明。
其中一份 199 页的稿件声称：$\zeta$ 与一切 Dirichlet $L$ 函数在半平面 $\Rea s>\tfrac78$ 中无零点（$s=1$ 处的极点除外），
并注明主要结果已在 Lean 中形式化，附有 comparator 检查的链接；稿件的后续应用部分未形式化。
另一份 49 页的稿件以相近方法给出 $\tfrac{11}{12}$ 的结论，未作形式化。
据 \emph{Quanta Magazine} 10 月 19 日的报道，专家大体认为 $\tfrac{11}{12}$ 的结果正确，已有人能用很短的篇幅概述其论证；
也有专家表示，在人工核对全部细节之前不视其为定理，并认为新方法未必使 RH 本身更近。上述稿件均未经同行评审。
本文未核查其 Lean 文件，以下只讨论这类声明应当怎样阅读。

若 $\zeta$ 在 $\Rea s>\theta$ 中无零点（$\theta<1$），由显式公式即得 $\psi(x)=x+O(x^{\theta}\log^2x)$；
这类结论统称拟黎曼猜想，此前没有任何 $\theta<1$ 被证明。依照前两节，对“已在 Lean 中形式化”的声明，读者可以依次追问：
\begin{enumerate}
 \item 挑战文件中的陈述是什么，是否就是稿件所声称的结论；例如例外点 $s=1$ 怎样处理，$L$ 函数用的是哪个定义。
 \item 证明依赖哪些公理，有无 \texttt{sorry}，是否经独立检查器复核。
 \item 形式化覆盖了稿件的哪些部分。
\end{enumerate}
前两项若都得到肯定的回答，结论的正确性就有了很强的保证，而且这种保证不依赖于有人读完 199 页。
但形式化回答不了另外的问题：证明的思想是什么，能否推广，是否新颖，结论归功于谁。
这些问题仍要靠人阅读和理解；上述报道中数学家的工作，主要也在于此。

\section{本书内容的形式化现状}
下表列出本书主线上若干结论在 mathlib 中的情况（以本仓库所用的 mathlib 版本为准，2025 年 8 月）。

\begin{center}\small
\begin{tabular}{@{}lll@{}}
\toprule
内容 & 书中位置 & 情况 \\
\midrule
Euler 乘积 & 第 5 章 & mathlib 已有 \\
$\Gamma$ 的倍元公式、余元公式 & 第 7 章 & mathlib 已有 \\
$\zeta$ 的延拓与函数方程 & 第 9 章 & mathlib 已有 \\
$\Rea s\ge1$ 上无零点 & 第 9 章 & mathlib 已有 \\
RH 的陈述；平凡零点 & 第 11 章 & mathlib 已有（RH 未证明） \\
$\xi$ 的 Hadamard 乘积 & 第 10 章 & 未收入 mathlib \\
零点计数公式 $N(T)$ & 第 11 章 & 未收入 mathlib \\
素数定理 & 第 12 章 & 未收入 mathlib；已在独立项目中完成 \\
显式公式 & 第 14 章 & 未收入 mathlib \\
RH 的算术等价形式 & 第 15 章 & 未收入 mathlib \\
\bottomrule
\end{tabular}
\end{center}

“未收入 mathlib”只说明还没有人把它形式化并收入这个库，不说明它不能形式化。从上节的例子看，
这些空白可以预期会由人与 AI 合作逐步补齐，那时本书的论证原则上也可以在形式上逐条核对；
而选择哪些陈述、怎样组织它们，仍然是数学写作本身的任务。

\section{结论}
\begin{enumerate}
 \item 形式化陈述必须给出每个对象的类型和定义；mathlib 让函数处处有定义，代价是在每条定理中写清例外点。
 $\zeta(1)$ 的约定值 $(\gamma-\log4\pi)/2$ 与函数方程的条件 \texttt{hs} 是两个具体例子。
 \item Lean 接受一个证明，保证的是在其逻辑和所列公理下由假设推出结论；它不保证陈述忠实于原意，
 假设矛盾的定理同样会被接受。
 \item 对不可信的证明者，可靠的做法是由可信的一方写定陈述（挑战文件），再用独立检查器核对证明。
 人工审查的对象于是从证明转移到陈述与定义。
 \item 形式化能确立正确性，不能代替理解。AI 时代的数学阅读，需要同时具备读懂形式化陈述和读懂数学思想两种能力。
\end{enumerate}

\section*{参考}
\begin{itemize}
 \item The Lean Language Reference, Validating Proofs. \url{https://lean-lang.org/doc/reference/latest/ValidatingProofs/}
 \item D.\ Loeffler, M.\ Stoll, Formalizing zeta and $L$-functions in Lean, \emph{Annals of Formalized Mathematics}, 2025. \url{https://arxiv.org/abs/2503.00959}
 \item PrimeNumberTheoremAnd. \url{https://github.com/AlexKontorovich/PrimeNumberTheoremAnd}
 \item J.\ D.\ Lichtman, Gauss --- an agentic formalization of the Prime Number Theorem, Fields Institute 报告，2025 年 10 月。 \url{https://www1.fields.utoronto.ca/talks/Gauss-agentic-formalization-Prime-Number-Theorem}
 \item 量子位，2025 年 9 月 14 日关于 Gauss 完成强形式素数定理形式化的报道。 \url{https://www.qbitai.com/2025/09/331038.html}
 \item J.\ Avigad, K.\ Donnelly, D.\ Gray, P.\ Raff, A formally verified proof of the prime number theorem, \emph{ACM Trans.\ Comput.\ Logic} 9 (2007).
 \item J.\ Harrison, Formalizing an analytic proof of the prime number theorem, \emph{J.\ Automated Reasoning} 43 (2009).
 \item AI proved the quasi-Riemann hypothesis. What happens now?, \emph{Quanta Magazine}, 2026 年 10 月 19 日。 \url{https://www.quantamagazine.org/ai-proved-the-quasi-riemann-hypothesis-what-happens-now-20261019/}
 \item 本仓库 \texttt{Lean示例/}；学习笔记第 10 篇（Li 判据）、第 16 篇（计算与证明）；本书第 9、11 章。
\end{itemize}
\end{document}
